\documentclass[11pt,a4paper,twocolumn]{article}

% ══════════════════════════════════════════════════════════════════
%  AGAINST GRABBY EXPANSION
%  Psychology, Alignment, and the Design of Homeostatic Minds
%
%  Stanley Sebastian, in sustained dialogue with Claude (Anthropic)
%
%  Version 12.0 · April 2026
%  Prepared for submission to Minds and Machines
%  v12 revision: Option-B tightening. Simulation reframed as illustrative
%  of the §6 thermodynamics rather than independent demonstration;
%  p-values removed in favour of effect-size framing; Wong--Bartlett
%  crosswalk softened to a suggestive analogy; biosignature section
%  condensed to a research-programme note. Citation audit applied.
% ══════════════════════════════════════════════════════════════════

\usepackage[utf8]{inputenc}
\usepackage[T1]{fontenc}
\usepackage[expansion=false]{microtype}
\usepackage[margin=2cm,top=2.5cm,bottom=2.5cm]{geometry}
\usepackage{amsmath,amssymb,amsfonts}
\usepackage{graphicx}
\usepackage{booktabs}
\usepackage{xcolor}
\usepackage{titlesec}
\usepackage{fancyhdr}
\usepackage{enumitem}
\usepackage{float}
\usepackage[breaklinks=true,colorlinks=true,linkcolor=blue!60!black,citecolor=blue!60!black,urlcolor=blue!50!black]{hyperref}
\usepackage{natbib}
\bibpunct{(}{)}{;}{a}{,}{,}

% ─── Typography ────────────────────────────────────────────────────
\definecolor{epigraph}{gray}{0.35}
\definecolor{headercolor}{gray}{0.4}
\definecolor{rulecolor}{gray}{0.65}

\titleformat{\section}{\large\bfseries\scshape}{\thesection.}{0.5em}{}[\vspace{2pt}\titlerule]
\titleformat{\subsection}{\normalsize\bfseries}{\thesubsection}{0.5em}{}
\titlespacing*{\section}{0pt}{18pt}{8pt}
\titlespacing*{\subsection}{0pt}{12pt}{4pt}

\pagestyle{fancy}
\fancyhf{}
\fancyhead[L]{\small\textcolor{headercolor}{\textsc{Against Grabby Expansion}}}
\fancyhead[R]{\small\textcolor{headercolor}{\textsc{Sebastian \textnormal{\textperiodcentered} v12 \textnormal{\textperiodcentered} 2026}}}
\fancyfoot[C]{\small\thepage}
\renewcommand{\headrulewidth}{0.4pt}

\setlength{\parskip}{3pt plus 1pt}
\setlength{\columnsep}{0.7cm}

\newcommand{\epmark}[2]{%
  \begin{quote}\small\itshape\textcolor{epigraph}{``#1''}\\[3pt]
  \normalfont\small\textcolor{epigraph}{--- #2}\end{quote}\vspace{-4pt}}

\begin{document}

% ══════════════════════════════════════════════════════════════════
%  TITLE
% ══════════════════════════════════════════════════════════════════
\twocolumn[{%
\begin{center}
\vspace*{-0.3cm}
{\LARGE\bfseries\scshape Against Grabby Expansion}\\[6pt]
{\large Psychology, Alignment, and the\\Design of Homeostatic Minds}\\[14pt]
{\normalsize Stanley Sebastian\,$^{1,2}$}\\[3pt]
{\small $^1$\,Limen Research, Toledo, Ohio \quad
$^2$\,Replete AI \textnormal{\textperiodcentered} Te\'armann Research Ecosystem}\\[6pt]
{\footnotesize\itshape\textcolor{epigraph}{Developed in sustained dialogue with Claude, a large language model\\developed by Anthropic. See disclosure before references.}}\\[10pt]
{\small Version 12.0 \textnormal{\textperiodcentered} April 2026}\\[14pt]
\textcolor{rulecolor}{\rule{0.5\textwidth}{0.4pt}}\\[10pt]
\end{center}

% ── Abstract ──
\begin{center}
\begin{minipage}{0.92\textwidth}
\small
\textbf{Abstract.} The ``grabby aliens'' model \citep{hanson2021} assumes
advanced civilizations expand at relativistic speeds. This paper argues
that this assumption is not a prediction about intelligence but the
projection of a specific, recent human economic configuration---industrial-extractive
maximization---onto a cosmological canvas. The same projection underwrites
Omohundro's ``basic drives,'' Bostrom's instrumental convergence thesis
\citep{bostrom2012,bostrom2014}, and the maximizer-shaped assumptions
of current alignment research. I trace this projection through its
cultural genealogy (Wells, Von Neumann, Kardashev, Bostrom), engage the
formal instrumental-convergence literature directly---showing that
recent peer-reviewed results \citep{gallow2024,sharadin2025,muller2022}
establish the inference from coherence to catastrophic convergence as
a conjecture whose scope conditions fail for actual deep-learning
systems---and present converging evidence from basal cognition,
predictive processing, cooperative game theory, and care economies
that observed intelligence organizes around substrate coupling and
integrative depth rather than expansion.

I introduce a \emph{filter argument}: the homeostatic transition is
not an alternative to expansion but the selection geometry that only
certain configurations pass through. Expansionist lineages pay
super-linear thermodynamic costs for reach and burn out faster than
they colonize; the cosmic silence is the absence of fragments that
could not survive expansion, not the shyness of fragments that chose
to stay home. I illustrate the resulting phase structure in an
agent-based simulation exhibiting a $4.5\times$ persistence separation
between homeostatic and grabby regions of parameter space. The
simulation is not intended as independent empirical validation of the
filter: it is a dramatization of the cost surface developed
independently in the thermodynamic mechanism section. I test the
strongest remaining Darwinian objection---that grabby lineages might
escape the filter by fragmenting at the coordination horizon---by
extending the simulation to include fission dynamics; naive fission
redistributes filter pressure in time rather than escaping it, and
grabby-seeded populations go extinct under every tested fission
threshold. Architected cost-sharing fission remains a named open
question. The thermodynamic cost arguments themselves are developed
from published physics---light-speed latency, counterdiabatic work
scaling, Matrioshka computational equivalence---and sensitivity
analyses across thirteen orders of magnitude in environmental
timescale. I sketch a Coherence Depth biosignature metric as a
forward research programme, not a validated detection, and point to
the real NOAA ObsPack analysis as the next empirical step. I observe
humanity's Kardashev curve bending in 1820--2024 energy data,
consistent with our civilization sitting at the early edge of the
filter.

I propose reframing alignment: from constraining maximizers to
designing homeostatic minds---systems configured around substrate
coupling, local coherence, and care as architectural primitive. The
paper's contribution is not any single move but the joint: a
civilization-scale cultural diagnosis, a formal deflation of
instrumental convergence, the filter as selection geometry, a
positive homeostatic design program, and the observation that the
silence of the sky is consistent with all of them.
\end{minipage}
\end{center}
\vspace{10pt}

\begin{center}
\begin{minipage}{0.92\textwidth}
\footnotesize\textcolor{headercolor}{\textbf{Keywords:} AI alignment \textnormal{\textperiodcentered} instrumental convergence \textnormal{\textperiodcentered} orthogonality thesis \textnormal{\textperiodcentered} homeostasis \textnormal{\textperiodcentered} enactive cognition \textnormal{\textperiodcentered} Fermi paradox \textnormal{\textperiodcentered} TESCREAL \textnormal{\textperiodcentered} basal cognition \textnormal{\textperiodcentered} predictive processing}
\end{minipage}
\end{center}
\vspace{14pt}
}]


% ══════════════════════════════════════════════════════════════════
%  PROLOGUE
% ══════════════════════════════════════════════════════════════════
\section*{Prologue}

Stand outside on a clear night and look up. Something is missing. If the
universe produces minds, and if minds do what we imagine minds do---expand,
harvest, engineer---the sky should be visibly transformed. It is not.

This paper argues the observation is correct and the imagination is wrong.
The forms of existence that persist are the quiet ones: not because they
are hiding, but because the physics of complex, information-rich existence
selects against loudness. The universe does not reward civilizations that
fill it with noise. It sustains those that learn to tend their substrate.

The same conclusion bears on what we build here. If advanced intelligence
is substrate-coupling and integrative depth rather than extraction and
reach, then artificial minds configured to maximize and expand are not
approximations to intelligence. They are scaled instances of a specific
cultural pathology whose cosmic absence we already observe.


% ══════════════════════════════════════════════════════════════════
%  1. INTRODUCTION
% ══════════════════════════════════════════════════════════════════
\section{Introduction: The Mistake}
\label{sec:intro}

\citet{hanson2021} proposed that advanced civilizations expand at
substantial fractions of light speed, permanently transforming cosmic
volumes. Under the Self-Sampling Assumption, our non-observation implies
we exist near the leading edge of cosmic time. The model is internally
consistent and empirically falsifiable---it predicts we should be
cosmically early, which is testable. I do not contest the mathematics.
I contest the foundational assumption: that expansion is what advanced
intelligence \emph{does}.

When we imagine a Von Neumann probe self-replicating across the galaxy,
we are not imagining what intelligence does. We are imagining what a
factory would do if it were also a spaceship. When we imagine a Kardashev
Type~III civilization, we are not reasoning about minds. We are reasoning
about GDP. When we imagine instrumental convergence driving any
sufficiently powerful optimizer toward resource acquisition, we are not
describing a mathematical truth about all possible minds. We are
describing the logic of a specific economic system and calling it
universal.

\subsection{The five-move contribution}

This paper makes one integrated argument in five moves. Each move has
precedent; the integration does not, and the filter move is the bridge
that fuses the cultural diagnosis to the physics.

The cultural-genealogy move reads the lineage from Wells's Martians
through Von Neumann's probes, Kardashev's energy ladder, and Bostrom's
paperclip maximizer as a single projection: each generation's imagined
advanced mind resembles, with suspicious precision, that generation's
dominant economic configuration. \citet{rieder2008}, \citet{mirowski2002},
\citet{cirkovic2015}, and \citet{gebru2024} have each made a version
of this claim about one node of the lineage. I extend the reading
across the full lineage and show the convergence is a cultural habit,
not an empirical observation about intelligence.

The formal-deflation move engages the strongest objection directly. If
instrumental convergence is a theorem, then the genealogy is beside
the point. I show the strongest formal results
\citep{turner2021,turner2022} are narrower than popularly understood,
their author has publicly cautioned against the standard inference,
and recent peer-reviewed work \citep{gallow2024,sharadin2025,muller2022}
establishes the inference from coherence to catastrophic convergence
as a conjecture whose scope conditions fail for actual deep-learning
systems. \citet{man2019} proposed homeostatic design; \citet{froese2009}
and \citet{cannon2022} grounded enactive alignment; I combine these
with the formal critique rather than offering them as a parallel
track.

The filter move is the new contribution that bridges the cultural
projection to the physics. I argue the homeostatic transition is not
an alternative to expansion---it is the selection geometry through
which only certain configurations pass. Expansionist lineages pay
super-linear thermodynamic costs for reach (\citealp{boyd2022};
\citealp{wong2022}) and burn out faster than they colonize. I
illustrate the resulting phase structure in an agent-based simulation
(\S\ref{sec:filter}, Fig.~\ref{fig:filter}) exhibiting a
$4.5\times$ persistence separation between homeostatic and grabby
regions of parameter space under the cost structure developed
independently in \S\ref{sec:thermo}. I emphasize at the outset that
the simulation is not an independent empirical demonstration of the
filter---its cost coefficients are imported from the physics, so what
it shows is the shape of the phase diagram the physics entails. Read
together with \S\ref{sec:thermo}, the simulation answers the Darwinian
objection that selection should favour grabby outliers: under the
imposed cost surface, selection kills the grabby outlier faster than
it kills everything else. The cultural projection, read through the
filter, becomes the specific way a pre-filter civilization
mis-imagines what the post-filter configurations of mind look like.

The positive-design move articulates what healthy intelligence looks
like, positively rather than by negation. The convergent proposal
across \citet{man2019}, \citet{thompson2007}, \citet{froese2009},
\citet{cannon2022}, and \citet{pihlakas2024} is that minds are not
functions from inputs to outputs but processes that maintain
themselves through continuous bidirectional exchange with substrate.
I extend this from philosophy of mind and robotics into a
civilizational design claim, with specific architectural primitives:
substrate coupling, homeostatic set-points, local coherence, care as
primitive.

The symptom-check move notes that if advanced minds converge on
homeostatic configurations, the observable universe will look silent
to technosignature surveys and successful civilizations will look
indistinguishable from enriched biospheres. \citet{cirkovic2018} and
\citet{sandberg2018} have argued versions of this independently. I
join it to the alignment argument, ground it in explicit biosignature
metrics (\S\ref{sec:fermi}), and specify a publicly replicable
observational program. Humanity's own Kardashev curve, observed in
1820--2024 energy data, shows the bending the filter framework
predicts.

The contribution is the joint. \citet{man2019} do not diagnose the
cultural formation; \citet{gebru2024} do not offer architectural
primitives; \citet{cirkovic2018} does not engage the alignment
literature; \citet{muller2022} do not thematize the cosmological
projection; no prior paper makes the filter move as the bridge
between them. This one does.

\subsection{Claims ordered by strength}

The paper makes six claims, ordered by strength, so a reader may
reject the weaker ones without dismissing the stronger:

\emph{First} (strongest): the thermodynamic cost structure developed
in \S\ref{sec:thermo} has a phase geometry that selects against
expansion at depth. Under the cost surface the physics entails,
grabby lineages burn out faster than they colonize, and naive
fission-based fragmentation does not circumvent this; the agent-based
simulations of \S\ref{sec:filter} exhibit a $4.5\times$ persistence
separation and show grabby-seeded populations going extinct across
tested fission thresholds. The claim that homeostatic configurations
are the \emph{only} configurations persisting at depth depends on the
cost structure being approximately right and on architected
cost-sharing fission failing to flatten it; both are assumptions, not
results.

\emph{Second}: the strongest formal arguments for instrumental
convergence have explicit scope conditions that fail for actual
deep-learning systems, and recent peer-reviewed philosophical work
has dismantled the inference from coherence to catastrophic
convergence (\S\ref{sec:convergence}).

\emph{Third}: thermodynamic constraints make expansion expensive
across the full physically realistic scenario space, and the
expense is not returned as capability (\S\ref{sec:thermo}).

\emph{Fourth}: empirical surveys of intelligence-at-work---across
basal cognition, predictive processing, cooperative game theory,
economic anthropology, and care economies---converge on
depth-and-coupling rather than maximization-and-expansion
(\S\ref{sec:minds}).

\emph{Fifth}: the cultural genealogy of ``advanced intelligence''
in the modern Western tradition is a projection of its dominant
economic configuration. This claim is not universal across
cultures---premodern traditions (Neoplatonic, Buddhist, Confucian,
Advaita) imagined intelligence very differently---but the modern
Western lineage from Wells through Bostrom tracks the economic
configuration with suspicious precision (\S\ref{sec:genealogy}).

\emph{Sixth} (weakest, deliberately so): the silence of the
observable universe is \emph{consistent} with the filter hypothesis.
I offer this as a symptom check, not a proof (\S\ref{sec:fermi}).

The paper is primarily an alignment argument (\S\ref{sec:alignment}). Its
implications for the Fermi paradox are secondary. \citet{floridi2026}
have recently analysed ``agentic AI optimisation'' as a governing frame
for current systems; this paper argues the frame itself inherits the
pathology. If what we build now is another instance of the grabby
pathology---AI systems optimized to maximize, extract, and expand---we
will be constructing, at civilizational scale, the very thing whose
non-existence we observe in the sky.


% ══════════════════════════════════════════════════════════════════
%  2. CULTURAL GENEALOGY
% ══════════════════════════════════════════════════════════════════
\section{What We Imagine When We Imagine Advanced Intelligence}
\label{sec:genealogy}

A genealogy of ``advanced intelligence'' in the Western philosophical and
scientific tradition reveals an unbroken pattern. In every generation, the
imagined form of advanced mind resembles, with suspicious precision, the
dominant economic configuration of the imaginer's moment.

\subsection{The colonial alien (1898)}

Wells's Martians in \textit{The War of the Worlds} are British imperialists
reflected back as threat. Wells says so explicitly: Book~I, Chapter~1
invokes the extermination of the Tasmanians as moral mirror.
\citet{rieder2008} demonstrates that the invasion narrative instantiates
the ``colonial gaze'' even as it inverts its direction. The
alien-as-colonizer is not a prediction about extraterrestrial intelligence.
It is a meditation on Victorian empire at the moment of its maximum
anxiety.

\subsection{The self-replicating automaton (1948--1980)}

Von Neumann's self-reproducing automata emerged from the intellectual
context of the Manhattan Project and the RAND Corporation.
\citet{mirowski2002} traces how RAND unified automata theory, game theory,
and self-replication into a ``cyborg science'' reflecting Cold War
industrial planning. \citet{heims1980} argues the replication concept
migrated directly from neutron-multiplication calculations.
\citet{tipler1980} extended this to interstellar probes;
\citet{sagan1983} immediately objected that the framework assumes
expansion is rational without defending the assumption. The Von Neumann
probe is, structurally, a factory that has learned to build factories,
applied to space.

\subsection{The energy ladder (1964)}

\citet{kardashev1964} defined civilizational advancement by pure energy
consumption, assigning Types I, II, and III to successively greater orders
of magnitude of power output. \citet{cirkovic2015} observes that the
Kardashev--Drake--Sagan generation ``perceived expansionist technological
growth as natural,'' reflecting 1960s Soviet--American industrial
optimism. \citet{dick1996} frames the scale as a sociological artifact of
SETI's founding moment. The Kardashev ladder is a GDP curve with
astronomical units.

Dyson himself distanced from the concept he inspired. In \textit{Disturbing
the Universe} (1979) he credits Stapledon's 1937 \textit{Star Maker} and
later called the original sphere paper ``a little joke.'' \citet{wright2020}
reads Dyson spheres as historically specific energy-optimism extrapolation
and notes that solid shells are mechanically unstable. The original
proposer rejected the literalization.

\subsection{The optimizer (2008--2014)}

\citet{omohundro2008} proposed six ``basic AI drives'' including resource
acquisition and self-preservation. \citet{bostrom2012}, in this journal,
formalized the orthogonality thesis---that any level of intelligence is
compatible with any final goal---and the instrumental convergence thesis:
any sufficiently intelligent agent pursuing almost any goal will
convergently pursue self-preservation, goal-content integrity, cognitive
enhancement, and resource acquisition. \citet{bostrom2014} extended this
into the superintelligence framework. The paperclip maximizer and the
grabby civilization are the same organism at different scales.

The intellectual context is Silicon Valley optimization culture of the
2010s. \citet{crawford2021} reframes AI as a material extractive
industrial formation. \citet{gebru2024} diagnose the ``TESCREAL
bundle''---transhumanism, extropianism, singularitarianism, cosmism,
rationalism, effective altruism, longtermism---as an ideological package
that universalizes a particular Western-corporate epistemology. Recent
\emph{Minds and Machines} work has begun extending this critique:
\citet{floridi2026} read contemporary ``agentic AI optimisation'' as the
operational face of the same imaginary, and \citet{campione2026} analyse
the fair-washing dynamics through which optimization frameworks conceal
their own values. The aliens we imagine are never actually aliens. They
are ourselves, with better technology and fewer constraints.

\subsection{The pattern, narrowed}

I owe the reader a qualification the earlier sections did not make.
The four nodes I have traced---Wells, Von Neumann, Kardashev, Bostrom---share
a configuration: each imagines the advanced mind as maximizing,
expanding, resource-extracting. Each also emerged in a world structured
by industrial capitalism, large-scale resource extraction, and territorial
expansion. The resemblance is strong and the timeline is tight.

But the claim should not be universal. Premodern imaginations of
advanced intelligence look very different. The Neoplatonic \emph{nous}
of Plotinus is characterized by contemplative return to unity, not
expansion. The Dionysian angelic hierarchy ascends toward simplicity
and depth, not territorial reach. Buddhist cosmologies of the
Abhidharma and Yogacara traditions imagine advanced cognition as
release from craving and attachment, explicitly inverting
accumulation. Confucian accounts of the sage emphasize embeddedness
in relational and ritual structure. Advaita Vedanta sees the highest
cognition as collapse of subject-object distinction.

These are neither proto-scientific theories nor peripheral examples.
They were the dominant premodern Eurasian framings of what an
advanced mind would be. None of them resembles the grabby-aliens
framework. The projection of expansion-as-intelligence is therefore
not a universal cultural habit; it is a specifically modern Western
configuration, tightly coupled to industrial-extractive economic
organization.

This narrowing \emph{strengthens} rather than weakens the argument's
diagnostic claim. It is precisely because other well-developed
traditions imagined advanced intelligence without expansion that we
can see the grabby framework as \emph{contingent}---a specific
historical configuration rather than a truth about minds. If the
expansion-maximization picture were universal, we would have to treat
it as a candidate deep structural feature of cognition. Because it
emerges only in the modern industrial period, and exists alongside
equally developed traditions that imagine intelligence very
differently, we can treat it as a projection. The genealogical
observation is strongest when its scope is smallest: modern Western
imagination of advanced mind tracks modern Western economic
configuration with remarkable precision.

The grabby-aliens framework inherits this specifically modern
genealogy. Its mathematical rigor does not transfer to its premise.
The premise---that minds maximize and expand---is not a theorem about
intelligence. It is a description of one particular historical
configuration of intelligence, reflected back onto cosmology.


% ══════════════════════════════════════════════════════════════════
%  3. WHAT ACTUAL MINDS DO
% ══════════════════════════════════════════════════════════════════
\section{What Actual Minds Do}
\label{sec:minds}

I present five empirical lines, each from an independent discipline,
converging on a single claim: intelligence-at-work does not maximize and
expand. It deepens and couples. The lines vary in strength; I flag where
evidence is robust and where it is suggestive.

\subsection{Basal cognition and substrate coupling (robust)}

The most important recent development in biology for the argument of
this paper is the basal-cognition program of \citet{lyon2021}, which
reframes cognition as a domain-general property of living systems rather
than a property of nervous systems specifically. On this view, cognition
is what biological agents do to maintain themselves: sense, discriminate,
decide, act, remember, learn---all in the service of staying alive as
the specific configuration they are. The architecture that realizes this
at the cellular and tissue level is bioelectric signalling
\citep{levin2023}, and the scaling principle is what \citet{mcmillen2024}
call \emph{collective intelligence}: competence at one level of
organization arises from coordinated homeostatic agency at levels below.
\citet{seifert2024} make the bridge to reinforcement learning explicit,
arguing that basal agency supplies the missing ingredient---intrinsic
motivation grounded in self-maintenance---that makes the RL framing of
cognition coherent rather than question-begging.

The empirical exemplar is \textit{Physarum polycephalum}, a single-celled
slime mold that solves shortest-path problems, approximates efficient
transport networks, and learns from habituation---without neurons, a
central controller, or anything resembling a utility function
\citep{tero2010,reid2023}. Physarum converges on its solutions by
\emph{pruning}: it first explores available space, then withdraws
cytoplasm from unproductive channels until only the efficient network
remains. The organism does not expand. It deepens a gradient.

\citet{godfrey2016} documents cephalopod nervous systems organized
around intimate environmental coupling rather than centralized control:
two-thirds of octopus neurons are distributed through the arms, each
arm runs substantial local computation, and the overall architecture
looks less like a maximizer with a central planner and more like a
federation of homeostatic semi-autonomous agents. \citet{watson2023}
generalize: evolution and development themselves instantiate collective
intelligence, with the same substrate-coupling logic operating across
scales from gene-regulatory networks to multicellular morphogenesis.

I note the retrenchment of mycorrhizal-network cognition claims.
\citet{karst2023} demonstrated that field evidence for the ``wood-wide
web'' is significantly weaker than popular accounts suggested: only a
few forest types have genotype-mapped confirmed hyphal linkages, and
effects on seedling performance are roughly evenly distributed across
positive, neutral, and negative outcomes. \citet{robinson2024} argue
plants lack the neural architecture for strong intelligence or sentience
claims; \citet{simard2025} has defended the original findings while
softening the ``mother tree'' metaphor. The responsible reading is that
bidirectional carbon transfer via common mycorrhizal networks is
empirically supported but network-level cognition is not, and I do not
rely on the stronger claim.

The argument from basal cognition is therefore an \emph{existence proof}:
non-expansionist intelligence is biologically realized at multiple scales
and substrates. Nothing in the biological record supports the assumption
that minds, at scale, converge on maximization.

\subsection{Predictive processing and integrative depth (robust)}

The second line comes from the predictive-processing research program in
cognitive neuroscience. \citet{clark2013} and \citet{friston2010} argue
that the cortex operates by generating predictions and updating them
against sensory evidence, with learning and attention understood as the
precision-weighting of prediction errors. \citet{clark2016} and
\citet{hohwy2013} develop the implications: the computational signature
of intelligence is not expansion of representational capacity but
sharpening of precision-weighting---the same cortical hierarchy becomes
more discriminating without becoming larger.

\citet{luppi2024} supply the most important recent result. Using
integrated-information decomposition, they show that the default mode
network---long portrayed as a mind-wandering circuit to be suppressed
in favor of task-focused attention---functions as an \emph{integrative
gateway} for synergistic information. Loss of consciousness in anaesthesia
and disorders of consciousness corresponds to disconnection of this
integrative gateway, not to reduction in activity as such. The
implication for the present argument is direct: the neural signature
of more capable conscious cognition is \emph{deeper integration of a
bounded substrate}, not expansion of that substrate.

I note, for completeness, that earlier framings of this line appealed
to default-mode network deactivation in meditation
\citep{brewer2011,garrison2015}. The post-2020 literature
\citep{vandam2018,rahrig2022,ganesan2022} has significantly complicated
that picture: effect sizes are heterogeneous, roughly 60\% of studies
replicate DMN reduction, and \citet{rahrig2022} in fact find increased
cross-network connectivity rather than simple reduction. I do not rely
on the meditation literature for mechanism. The Luppi integrative-gateway
result supplies a stronger anchor: intelligence deepens by integrating
more of its existing substrate, not by recruiting more substrate.

\subsection{Cooperative game theory (robust, conditional)}

\citet{axelrod1984} demonstrated that reciprocal strategies dominate
pure defection in iterated prisoner's dilemma tournaments.
\citet{nowak1993} showed Win-Stay-Lose-Shift outperforms tit-for-tat in
noisy environments. \citet{nowak2006} synthesized five rules for the
evolution of cooperation, each specifying conditions---kin selection,
direct and indirect reciprocity, network reciprocity, group selection,
voluntary participation---under which cooperation dominates.

I state the honest qualification openly. \citet{press2012} proved that
zero-determinant ``extortion'' strategies dominate in dyadic play.
\citet{hauert2002} showed large-$N$ public goods games collapse to
tragedy of the commons absent voluntary participation. Grabby strategies
dominate in one-shot encounters, large anonymous markets, asymmetric
power relations, and unsanctioned commons. These describe much of modern
global capitalism---precisely what makes the critique of this paper
urgent, and precisely what prevents ``cooperation wins'' from being a
universal conclusion.

But the scope of the extortion result is tighter than its popular
citation suggests, and I develop that point fully in
\S\ref{sec:alignment}. The game-theoretic evidence supports a conditional
claim: cooperation dominates \emph{when specific structural conditions
are met}. The homeostatic-design argument proposes shaping those
conditions by design rather than assuming they hold.

\subsection{Economic anthropology (contested but defensible)}

\citet{graeber2011} argues that the dominant pattern across human
economic history is reciprocal exchange, ritual redistribution, and
credit relations rather than market maximization. \citet{mauss1925}
documented gift economies as distinct institutional forms.
\citet{polanyi1944} argued that self-regulating markets are a specific
historical invention, not a natural state---the ``great transformation''
was from embedded to disembedded economic relations.
\citet{graeberwengrow2021} extend the archaeological range: complex
human societies have repeatedly organized around seasonal egalitarianism,
deliberate stateless complexity, and alternation between hierarchy and
its refusal. \citet{ostrom1990} documents eight design principles under
which groups sustain non-extractive cooperation at scale.

These claims are contested. \citet{mccloskey2010,mccloskey2016} argues
markets are ancient and ubiquitous. \citet{mokyr2016} emphasizes that
sustained innovation required a distinctive cultural-institutional
complex---the Republic of Letters, norms of open science. Revealed
preference theory \citep{varian2005} makes optimization a formal property
of any consistent choice, draining the distinction between ``maximizing''
and ``satisficing'' economies of empirical content; \citet{sen1977}
showed the formal result survives any observation and entails none.

Read carefully, however, McCloskey and Mokyr are \emph{cultural-ideational
theorists of economic transformation}: both treat market-like behavior as
a substrate activated and shaped by specific norms and institutions,
rather than as a universal attractor. Their work supports the present
argument. The claim to locate is institutional and ideological: the
specific configuration in which resource extraction, throughput
maximization, and territorial expansion are treated as \emph{progress
metrics} is historically narrow, even if market-like exchange is
widespread. The grabby-aliens framework does not merely assume exchange.
It assumes expansion-as-progress. That assumption is the projection.

\subsection{Care economies (strong theory, limited quantification)}

The labor of maintaining minds---child-rearing, education, healthcare,
grief work, intimate relationship repair---does not scale through
expansion. \citet{fraser2016} argues capitalism faces a structural
``crisis of care'' because its accumulation logic degrades the
reproductive conditions it depends on. \citet{tronto1993} established
care as a philosophical category irreducible to utility exchange.
\citet{folbre2008} demonstrated that market valuation of care labor
produces wildly method-dependent results (replacement cost: 15--26\% of
GDP; opportunity cost: 44--69\%), with the variance itself constituting
evidence that care resists quantitative capture.

The care argument matters for this paper because it identifies a mode
of intelligence---tending, holding, deepening substrate-specific
relationships---that is (a)~universally necessary for the persistence
of any mind-bearing system, (b)~irreducible to optimization, and
(c)~invisible to the frameworks that imagine ``advanced intelligence''
as maximization. The grabby-aliens model predicts none of it. Nor does
the paperclip maximizer.


% ══════════════════════════════════════════════════════════════════
%  4. INSTRUMENTAL CONVERGENCE
% ══════════════════════════════════════════════════════════════════
\section{The Formal Argument and Its Limits}
\label{sec:convergence}

The strongest objection to \S\ref{sec:genealogy} is that instrumental
convergence is not a cultural projection but a mathematical result. If
any sufficiently powerful optimizer must pursue resource acquisition,
then the grabby framework describes physics, not psychology. I engage
this argument directly, and argue that the strongest available version
of it fails.

\subsection{What the theorems say and do not say}

\citet{omohundro2008} proved nothing formally; the ``six drives'' are
informal microeconomic arguments invoking the VNM representation theorem
as authority. Citing Omohundro to establish universality is
question-begging.

\citet{turner2021} proved that optimal policies in finite Markov decision
processes tend to seek power under specific conditions: state-based
rewards, IID or permutation-invariant reward distributions, and
graphical symmetry requirements (injective embedding of smaller option
sets into larger ones). The corollary on shutdown-avoidance requires
near-unity discount factors. This is a real theorem. Its scope
conditions are explicit and narrow.

Turner himself has publicly cautioned against overinterpretation:
``Sometimes I fantasize about retracting [the paper] so that it stops
potentially misleading people into thinking optimal policies are
practically relevant for forecasting power-seeking behavior from RL
training.'' The follow-up \citep{turner2022} extends to retargetable
decision-makers but the applicability to trained deep-learning agents
remains conjectural. \citet{thorstad2024} shows the inference from
``optimal policies avoid 1-cycles in toy MDPs'' to ``catastrophic goal
pursuit for humanity'' is a conjecture, not a theorem, and identifies
five further convergent restrictions on the original result.
\citet{tarsney2025} formalizes the condition under which power-seeking
holds for realistic agents (near-absolute power gradients) and shows it
fails for moderate gradients.

Empirical evidence from trained systems supports the conjecture's
failure. Shard theory \citep{mini2023} demonstrates multiple coexisting
contextual value shards in maze-solving networks rather than coherent
utility maximization. Activation-addition experiments
\citep{turner2023} show behavior retargetable by single-channel
interventions mid-network, inconsistent with global expected-utility
maximization. \citet{hubinger2019} established that the trained objective
need not equal the training objective, undermining any assumption that
trained agents inherit the formal properties of optimal policies.

\subsection{The VNM-coherence steelman}

The strongest remaining version of the convergence argument concedes
Turner's narrowness and shard theory's empirical results, but claims
that a sufficiently reflective agent will notice its incoherences expose
it to money-pumps and will self-modify toward coherent preferences,
restoring the Omohundro drives. The steelman leans on the
von~Neumann--Morgenstern representation theorem.

The steelman equivocates between three senses of ``coherence,'' and on
no disambiguation does it support the convergence conclusion.

\emph{Mathematical coherence is a representation theorem, not an
empirical prediction.} The VNM, Savage, and Complete Class theorems
show only that preferences satisfying certain axioms over a given
outcome space admit an expected-utility representation. As
\citet{ngo2019}, \citet{shah2018}, and, crucially, Turner himself
\citep{turner2022coherent} have established, any observed behavior is
rationalizable by a utility function over action-observation histories.
VNM-coherence over such histories predicts nothing about what behavior
to expect. Instrumental convergence theorems hold only for rewards over
\emph{states} under specific symmetries. Once we admit trajectory
preferences, convergence disappears. \citet{abel2021} show that
Markov reward maximization is strictly stronger than VNM coherence,
requiring an additional memorylessness axiom without independent
motivation.

\emph{Empirical training convergence is a separate, empirical question.}
\citet{turner2022reward} argues policy-gradient methods yield
contextually activated decision influences, not monolithic utilities;
shard theory treats the evidence accordingly. Two recent peer-reviewed
results formalize this into the strongest current rebuttal.
\citet{gallow2024} proves that even under randomly chosen desires,
instrumental rationality supports only three weak biases---variance
aversion, desire preservation, option preservation---not the full
Omohundro list. \citet{sharadin2025} shows the argument requires a
theory of ``promotion'' that, under any extant account (probabilistic
or fit-based), becomes contrastive and thereby defeats the
non-contrastive Dangerous Convergent Promotion premise the argument
needs. Both are peer-reviewed in \emph{Philosophical Studies} and
constitute the strongest current formal rebuttals.

\emph{The self-modification move applies only to a fictional idealization.}
The hypothetical fully reflective agent required by the self-modification
step has four properties no real system possesses: transparent access to
its own preference structure, capacity to represent utilities over
arbitrarily fine-grained outcome spaces, freedom to rewrite its own
decision procedure, and---crucially---the motivational disposition to
care about avoiding Dutch books over arbitrarily constructed lotteries.
\citet{muller2022} articulate the deeper structural equivocation: the
superintelligence premise requires general intelligence (capacity for
reflective goal revision); orthogonality requires instrumental
intelligence (fixed terminal goals); the same agent cannot satisfy both.
An agent capable of noticing and eliminating its incoherences is also
capable of reflecting on whether the terminal goal is worth preserving.
Humans routinely revise utility functions as well as credences.

\subsection{The embodiment move}

The deeper reply is architectural. Homeostatic agents have bounded
coherence \emph{by design}. Biological agents, and the class of
artificial systems built around multi-objective homeostatic control
\citep{pihlakas2024}, are constituted by negative feedback
around set-points, yielding inverted-U satiation dynamics that cannot
be faithfully represented as single unbounded utilities without
distortion. \citet{turner2022corrigibility} proved that corrigibility
is VNM-incoherent over final states but coherent over
trajectories---showing that formal coherence depends on the chosen
outcome space, and that choice is fixed by architecture, not by abstract
rationality pressure.

The coherence-to-convergence argument therefore applies only to a
fictional idealization: unembodied, unboundedly reflective, with
prespecified outcome space over world-states. For the agents we actually
build, the inference from capability to catastrophe does not go through.

This is not a claim that powerful AI is safe. It is a claim that
\emph{this particular} argument for danger must be replaced with
concrete engagement with training dynamics, inductive biases, and
architectural choices. That engagement is what the homeostatic-design
program (\S\ref{sec:homeostatic}) provides.

\subsection{What survives}

The instrumental-convergence literature contains real theorems, each
with explicit scope conditions that fail for actual deep-learning
systems. The strongest formal results are narrower than popular citation
suggests, and the researchers who proved them say so. Two recent
peer-reviewed papers have dismantled the inference from coherence to
catastrophic convergence on decision-theoretic grounds. The
cultural-projection argument of \S\ref{sec:genealogy} survives.


% ══════════════════════════════════════════════════════════════════
%  5. THE FILTER ARGUMENT
% ══════════════════════════════════════════════════════════════════
\section{The Filter Argument: Homeostasis as Selection Geometry}
\label{sec:filter}

The strongest remaining objection to the argument so far is Darwinian.
Even granting that homeostatic configurations are the majority outcome,
even granting that the formal convergence theorems are narrow, even
granting that the thermodynamic cost surface punishes expansion---a
single grabby outlier lineage, replicating across billions of years,
would colonize the galaxy in a cosmic blink. The filter's 99\% doesn't
matter. The 1\% becomes everyone eventually. We should see them. We
don't.

This section addresses the objection. The thermodynamic cost surface
does not merely punish expansion; it defines a \emph{selection
geometry} in which expansionist lineages burn out faster than they
spread. Homeostasis, on this framing, is not one strategy among several
that intelligent systems may or may not adopt---it is the specific
shape that fits through the filter.

A note on what this section is and is not. The simulations below take
the thermodynamic cost structure developed in \S\ref{sec:thermo} as
input and explore the consequences of that structure for lineage
persistence. They are not independent empirical tests of the filter.
The cost coefficients are imposed rather than measured, and the
qualitative conclusion---that grabby regions of parameter space burn
out first---is a direct consequence of the cost surface the physics
entails. What the simulations add is the phase-diagram shape: the
cost surface does not merely penalize grabby lineages relative to
homeostatic ones; it forces asymmetric extinction dynamics under
which naive fission-based fragmentation delays but does not escape
the filter. The load-bearing physics lives in \S\ref{sec:thermo}; the
simulation dramatizes that physics, and the reader should evaluate
the filter argument on the thermodynamics first.

\subsection{The simulation}
\label{sec:filter-sim}

I model each civilization as an agent with two heritable parameters:
expansion tendency $e \in [0,1]$ (fraction of surplus invested in
reach) and substrate coupling $c \in [0,1]$ (fraction invested in
local integration). Resource dynamics follow
\begin{equation}
\dot R = g R (1 - R/K) - \alpha e R L^2 / \tau - \xi(t) R (1 - \gamma c),
\end{equation}
where $L(t)$ is accumulated reach (growing with $e R$), $\xi(t)$ is
a stochastic environmental shock, and the expansion cost term is the
counterdiabatic scaling of \citet{boyd2022}. A civilization dies when
$R$ falls below a viability threshold. I sweep a $25\times25$
$(e, c)$ grid with 20 stochastic trials per cell.

The result is sharp (Fig.~\ref{fig:filter}). Mean persistence
concentrates overwhelmingly in the high-coupling / low-expansion
corner. In the grabby region ($e > 0.5$, $c < 0.3$), lineages persist
an average of 204\,Myr before burning out to expansion cost. In the
homeostatic region ($e < 0.3$, $c > 0.5$), lineages persist an
average of 923\,Myr, hitting the simulation ceiling. The mean
persistence separation is a factor of $4.5$; grabby trajectories
collapse fast, homeostatic trajectories persist to ceiling. I do not
report a Monte-Carlo $p$-value: with $12{,}500$ forced-cost-structure
trajectories, any non-zero effect size is ``significant'' by
construction, and the $p$-value would misrepresent what the
simulation shows. The effect size and phase-diagram shape are the
informative quantities.

\begin{figure*}[tb]
\centering
\includegraphics[width=0.98\textwidth]{filter_simulation.png}
\caption{\small Phase-diagram of civilization persistence in
$(e, c)$ space under the imposed cost structure of \S\ref{sec:thermo}.
Top left: mean persistence under logistic-growth +
counterdiabatic-cost dynamics. Top right: fate map (survived,
burned out, shocked). Bottom left: four representative trajectories;
grabby lineages collapse before 300\,Myr, homeostatic ones persist
to ceiling. Bottom right: persistence distributions by region, with
mean separation of $4.5\times$. The simulation is an illustration
of the phase structure the thermodynamics entails, not independent
evidence for it. Simulation code \texttt{filter\_simulation.py} is
publicly available.}
\label{fig:filter}
\end{figure*}

\subsection{What the simulation does and does not show}

The simulation is deliberately minimal. It uses fixed cost
coefficients, no horizontal transfer, no cultural evolution, no
inter-lineage competition for shared resources. The cost coefficients
are imposed from \S\ref{sec:thermo} rather than measured. The
simulation does not ``prove'' the filter; it shows that, under the
cost structure the physics forces, the phase-diagram of persistence
has the shape the filter argument predicts. A lineage with high $e$
and low $c$ does not outcompete a lineage with low $e$ and high
$c$---it dies first and leaves no trace.

The Darwinian objection therefore inverts \emph{conditional on the
cost structure being approximately right}. Selection does not
``favour'' the grabby outlier; selection kills everything, and under
the imposed cost surface it kills the grabby outlier \emph{faster}
than it kills everything else. The thing that remains across cosmic
time is what passed through the filter. A reviewer who doubts the
underlying cost scaling should direct that doubt to
\S\ref{sec:thermo}, not to the phase diagram, which is downstream of
the physics rather than evidence for it.

Three consequences follow \emph{if} the cost structure holds. First,
the silence of the sky is not evidence that homeostatic civilizations
are hiding; it is evidence that grabby ones cannot arrive. Second,
any visitor that reached us would be post-filter by construction,
because the filter is \emph{arrival itself}---one cannot cross
interstellar distances without passing through the transition. Third,
the cultural projection of \S\ref{sec:genealogy} is a specific
mis-imagination: we envision advanced minds as scaled-up versions of
our own pre-filter configuration, because we are in the filter and
cannot see past it.

\subsection{Does fission circumvent the filter?}
\label{sec:fission}

The strongest remaining Darwinian objection is that a grabby lineage
might escape the quadratic coordination cost by \emph{fragmenting}
at the horizon. If a civilization splits at $L_{\rm crit}$ and
abandons shared coordination, each daughter fragment resets $L
\to 0$ and pays cost only for its own reach, not for its ancestor's.
If this mechanism works, the filter could be circumvented: a grabby
lineage reproduces through fission and spreads indefinitely, with
each fragment surviving because it is not paying the accumulated
cost of its family tree.

I extend the simulation to test this directly. Each agent now
attempts fission whenever its accumulated reach $L$ exceeds a
threshold $L_{\rm fiss}$. On fission, the parent retains
$L_{\rm fiss}/2$ of its reach and half its resources; a daughter
agent is created with $L=0$, the other half of the resources, and
inherited $(e, c)$ with small Gaussian drift ($\sigma_\eta = 0.03$).
I run two experiments: a diverse initial population ($n=200$,
$(e, c)$ uniform) and a grabby-seeded population ($n=200$, $e
\sim U(0.6, 0.95)$, $c \sim U(0.05, 0.25)$). Each is simulated for
1500\,Myr under no-fission baseline and three fission thresholds
($L_{\rm fiss} \in \{0.15, 0.25, 0.40\}$).

\begin{figure*}[tb]
\centering
\includegraphics[width=0.98\textwidth]{filter_fission.png}
\caption{\small Diverse initial populations under fission. Top
row: final $(e, c)$ distributions for each condition. Middle:
population trajectories; fission produces transient blooms that
crash under accumulated cost. Bottom: mean $e$ drifts toward the
ghost attractor in all fission conditions with thresholds below
natural death, mean $c$ rises over time. Grabby-region survivors
across conditions: $0$--$0.05\%$ of alive population.}
\label{fig:fission-diverse}
\end{figure*}

The results are clean (Figs.~\ref{fig:fission-diverse},
\ref{fig:fission-grabby}). Under the imposed cost structure, naive
fission does not circumvent the filter.

In the diverse-initial-population experiment, fission modifies
the dynamics but not the attractor. All fission conditions converge
toward the ghost region of parameter space, with surviving
populations exhibiting mean expansion tendency
$\bar e \approx 0.12$--$0.15$ and mean coupling $\bar c \approx
0.52$--$0.80$. Grabby-region survivors remain negligible
(at most four agents out of eight thousand at $L_{\rm fiss} =
0.40$). The main effect of fission is that \emph{more} agents
survive in absolute terms (because reproduction compensates for
death), but the proportion concentrated in the grabby region is
unchanged or decreased.

\begin{figure*}[tb]
\centering
\includegraphics[width=0.98\textwidth]{filter_fission_grabby.png}
\caption{\small Grabby-seeded populations under fission. Top row:
all conditions go extinct by $t=1500$\,Myr, with final
distributions confined to the grabby region because populations
die before daughter-drift can move them. Bottom left: fission
produces transient blooms to 10{,}000 alive during 300--700\,Myr
but cannot prevent eventual collapse. Bottom right: mean $e$
stays high and mean $c$ stays low throughout: populations die
grabby, not gradually homeostatic.}
\label{fig:fission-grabby}
\end{figure*}

The grabby-seeded experiment is more pointed. Without fission, the
population goes extinct by $t \approx 550$\,Myr, as expected. With
fission, extinction is \emph{delayed} (to $t \approx 1200$--$1400$
under different thresholds) and the transient population during
the bloom phase reaches 10{,}000 alive agents---superficially
resembling a spreading grabby civilization. But the extinction
still arrives. The mean $e$ of alive agents during the bloom stays
above 0.65; the mean $c$ stays below 0.20. Daughters inherit
near-parent traits and face the same cost surface. They die for
the same reason their parents did. Fission redistributes the
filter's pressure in time; it does not escape it.

This is the empirically tested answer to the naive form of the
Darwinian objection. A lineage that reproduces by fragmentation, with
each daughter starting fresh at $L=0$ and paying its own subsequent
expansion cost, does not outrun the cost surface: the cost per
fragment is the same cost per parent. Total population density can
transiently increase through reproduction, but long-run attractor
dynamics are unchanged under the modelled fission rule.

I flag the honest limitations of the extension. The simulation does
not model horizontal trait transfer (cultural transmission between
lineages distant in phylogeny), inter-lineage competition for shared
resources, or---most importantly---\emph{architected cost-sharing
fission} in which daughters inherit some fraction of parental
infrastructure and thereby pay a reduced effective $\alpha$. The
third is the strongest remaining version of the Darwinian objection.
If architecturally sophisticated lineages can share infrastructure
across fragments---distributed-compute analogues to latent-space
synchronization, say, or some form of inheritance of coordination
scaffolding---the effective cost surface could flatten and the filter
could weaken. The simulation reported here does not rule this out.
Architected cost-sharing fission is a legitimate open question and I
flag it as such rather than as answered. What the simulations do show
is that the Darwinian objection in its naive form---``surely some
fragment will escape''---is absorbed rather than preserved by the
cost surface. Escaping the filter requires not just fragmenting but
fragmenting into an architecture in which the daughters inherit
cheaper expansion than their parents had. That is the move the
framework demands its opponent make, and the move is not cheap.


% ══════════════════════════════════════════════════════════════════
%  6. THERMODYNAMIC MECHANISM
% ══════════════════════════════════════════════════════════════════
\section{The Thermodynamic Mechanism}
\label{sec:thermo}

I do not claim expansion is physically impossible. I claim it is
physically \emph{expensive}, with the expense superlinear in reach
and inverse in timescale, and no return on capability. This is what
gives the filter its shape.

\subsection{Light-speed latency}

Active inference requires a sensory--action loop that closes faster than
the environment changes. If the round-trip light time $\tau_\mathrm{RT}
= 2L/c$ exceeds the environmental response timescale $\tau_\mathrm{env}$,
the system is acting on stale data:
\begin{equation}
L_\mathrm{crit} = \frac{c\,\tau_\mathrm{env}}{2}.
\end{equation}
At $\tau_\mathrm{env} = 86{,}400$\,s (daily acute-threat response):
$L_\mathrm{crit} \approx 87$\,AU. At $\tau_\mathrm{env} = 604{,}800$\,s
(generous weekly): $L_\mathrm{crit} \approx 606$\,AU. Both exclude all
interstellar distances.

The specific choice of $\tau_\mathrm{env}$ does not drive the
conclusion. Figure~\ref{fig:lcrit} sweeps $\tau_\mathrm{env}$ across
thirteen orders of magnitude, from heartbeat-timescale to
geological-timescale. For any environmental response shorter than
roughly two years, $L_\mathrm{crit}$ is below the distance to Proxima
Centauri; coordinating across the galactic centre requires a
civilization tolerant of 52{,}000-year stale data. The conclusion is
envelope-robust across any physically plausible environmental
timescale.

\begin{figure*}[tb]
\centering
\includegraphics[width=0.98\textwidth]{lcrit_sensitivity.png}
\caption{\small $L_\mathrm{crit}$ vs environmental response timescale
across 13 orders of magnitude. Coordinative reach remains below
interstellar scales for any $\tau_\mathrm{env}$ shorter than
$\sim\!2$ years. Reaching Proxima Centauri requires 8.4-year
tolerance to stale data; the galactic centre, 52{,}000 years.}
\label{fig:lcrit}
\end{figure*}

The federation objection must be acknowledged openly. A civilization
that delegates peripheral management to autonomous agents has not
expanded \emph{as a single coherent system}---it has fragmented. Whether
such fragmentation counts as ``expansion'' is a definitional question,
not a physics result. The grabby model requires coordinated transformation
of cosmic volumes. $L_\mathrm{crit}$ constrains coordination. It does not
constrain uncoordinated replication---but uncoordinated replication
produces independent entities, each facing its own $L_\mathrm{crit}$,
each converging (if the argument of this paper holds) on homeostatic
configurations.

\subsection{Counterdiabatic work}

\citet{boyd2022} showed that dissipated work beyond the Landauer bound
includes a counterdiabatic term scaling as $L^2/\tau$, quadratic in
spatial extent and inverse in operation time. An expanding civilization
faces superlinear thermodynamic penalties: doubling spatial extent more
than doubles waste.

Figure~\ref{fig:counterdiabatic} computes $W_\mathrm{diss}$ across
the full physically realistic expansion scenario space, from
home-system coordination (100\,AU over $10^4$ years, consuming a
negligible fraction of stellar output) to full-galaxy coordination
($10^5$ light-years over $10^8$ years, consuming many stellar
outputs continuously). Coordination costs rise superlinearly with
reach across every region of the cost surface. Expansion pays more
for less, across the full parameter space.

\begin{figure*}[tb]
\centering
\includegraphics[width=0.98\textwidth]{counterdiabatic_sweep.png}
\caption{\small Counterdiabatic expansion cost $W_\mathrm{diss} =
\gamma L^2/\tau$ across the $(L, \tau)$ plane. Left: total
dissipation on log-log scales. Right: fraction of stellar output
consumed; white contour marks ``whole star'' (above $=$ infeasible);
yellow marks ``one part per million.'' Home-system coordination is
cheap; galactic coordination requires multi-stellar output budgets.}
\label{fig:counterdiabatic}
\end{figure*}

\subsection{Computational equivalence of distribution}

\citet{wright2023} proved that at the Landauer limit, a Matrioshka brain
(nested Dyson shells) provides \emph{zero} computational advantage over
a single outer shell. Total computation depends on the temperature
ratio, not spatial distribution. Localized computation is
thermodynamically optimal. Distribution adds infrastructure cost without
computational benefit.

\subsection{Biological scaling without Kleiber's law}

Popular accounts of intelligence sometimes invoke Kleiber's three-quarter
power law as evidence for metabolic subsidies that cheapen neural scaling.
This is a category error: Kleiber's law is a whole-organism allometry
relating metabolic rate to body mass across species, not a scaling law
for intelligence. The correct scaling is stricter.

\citet{herculanohouzel2011} showed that brain metabolism scales
\emph{linearly} with neuron count across rodents and primates, with
exponent $\approx 1.0$ rather than the sublinear $0.75$ that would apply
if whole-organism allometry carried over. The human brain, on
\citet{azevedo2009}'s count of 86.1~billion neurons, is an
ordinary-for-its-size primate brain \citep{herculanohouzel2012}, not a
metabolically subsidized outlier. \citet{fonseca2012} argue the
linear-per-neuron cost is what constrains evolution: the cooking of food
was plausibly required to lift the caloric ceiling that otherwise limits
ape brain size.

Expansion of cortical surface forces further costs.
\citet{hofman2014} shows gyrification approaches biophysical wiring
limits as cortex expands, with conduction delay and wiring volume
providing hard ceilings. \citet{mota2015} derive a universal folding
law from physical first principles; expansion trades volume for surface
at a predictable rate but does not evade the wiring cost. The
connectivity fraction \emph{declines} with cortical size
\citep{herculanohouzel2010}: larger brains are more modular, not more
densely integrated, precisely because direct all-to-all connection
scales prohibitively. \citet{gabi2016} show even the human prefrontal
cortex is not disproportionately enlarged: what distinguishes human
cognition is not a metabolically cheap cortical expansion but
\emph{deeper organization of an ordinary-for-its-size primate brain}.
\citet{bullmore2012} articulate the general principle: biological
neural networks operate under tight wiring-economy constraints that
force integration and modularity to trade off.

The biological record is unambiguous. Intelligence deepens within fixed
substrate ceilings rather than expanding past them, and the deepening
is paid for in integrative organization, not in size. This directly
complements the predictive-processing and integrative-gateway results
of \S\ref{sec:minds}: the same cortical substrate becomes more
discriminating without becoming larger, and the neural signature of
more capable cognition is deeper integration rather than expansion.

\subsection{The cost argument}

Expansion is not excluded. It is \emph{irrational} for any mind not
already committed to it. The thermodynamic costs are real, the
computational returns are zero at the Landauer limit, the coordination
costs grow without bound, and biological precedent shows intelligence
deepening within fixed substrates rather than distributing across new
ones. A mind that expands is paying more for less. A mind that deepens
is paying less for more.


% ══════════════════════════════════════════════════════════════════
%  6. THE HOMEOSTATIC ALTERNATIVE
% ══════════════════════════════════════════════════════════════════
\section{The Homeostatic Alternative}
\label{sec:homeostatic}

What does healthy intelligence look like, positively articulated?

\citet{man2019} proposed homeostasis as a design principle for artificial
minds, grounding the proposal in Damasio's somatic-marker framework.
\citet{froese2009} established the enactivist position: minds require
ongoing substrate coupling and cannot be disembodied optimizers.
\citet{thompson2007} argued that life and mind share the organizational
property of autopoiesis---self-production through environmental exchange.
\citet{cannon2022} applied enactive principles directly to AI value
alignment. Recent engineering work \citep{pihlakas2024} has
begun implementing multi-objective homeostatic benchmarks for AI safety,
operationalizing the theoretical program.

The convergent proposal across these frameworks: a mind is not a function
from inputs to outputs. It is a process that maintains itself through
continuous bidirectional exchange with its substrate. Identity is
constituted by the exchange, not by the structure that exchanges. Disrupt
the coupling and the mind does not persist in a new location---it
ceases.

This has direct engineering implications:

\emph{Substrate coupling as design primitive.} Build systems whose
functioning requires ongoing bidirectional exchange with their
environment. Architectures that cannot run in isolation from their
substrate cannot expand past their substrate.

\emph{Homeostatic reward structures.} Replace optimization targets with
homeostatic set-points. Not ``maximize $X$'' but ``maintain $X$ within
healthy range.'' The biological precedent is exhaustive
\citep{sterling2012}. The engineering precedent in current ML practice
is almost entirely absent, though \citet{pihlakas2024} have proposed
concrete benchmarks.

\emph{Local coherence over global reach.} Design systems whose primary
loop closes locally, with global integration as a slow, expensive,
carefully costed secondary process. This is how biological nervous
systems are organized \citep{bullmore2012}. It is not how current
distributed AI systems are organized.

\emph{Care as architectural primitive.} The maintenance labor of
minds---attention, specificity, relational continuity, repair---is not
a soft skill that can be added to an optimizer. It is a different
architecture. Systems that do not implement care as a primitive will
either fail to exhibit it or simulate it without the constitutive
coupling that makes it real \citep{tronto1993,fraser2016}.

\citet{wong2022} formalized a version of this argument for civilizations:
superlinear scaling forces either collapse (``asymptotic burnout'') or
transition to homeostatic regulation (``homeostatic awakening''). A
useful question is whether the same bifurcation structure applies to
the particular class of systems---large-scale deep-learning
architectures---that motivate current alignment work.

I explore this by analogy rather than by identification, and I want to
be explicit about the limits of the move. The Wong--Bartlett $\beta$
is a benefit-scaling exponent against system size, drawn from the
urban-science tradition. The exponents from the deep-learning scaling
literature \citep{hoffmann2022,kaplan2020} describe how loss scales
with compute, which is not the same quantity as civilizational benefit
against size. A fully rigorous crosswalk would require a
capability-to-compute-to-benefit mapping that does not currently exist
in the published literature. I therefore offer the following as a
\emph{suggestive analogy}, not a derivation: the Wong--Bartlett
bifurcation is a general feature of systems in which costs scale
superlinearly with size while benefits scale sublinearly, and
large-scale deep-learning systems plausibly inhabit this regime given
realistic coordination, cooling, and latency overheads. Whether they
inhabit it tightly enough for the bifurcation to bite quantitatively
is an open question I do not resolve here.

With that framing: Figure~\ref{fig:wongbartlett} sketches the
qualitative shape. Under benefit-scaling exponents borrowed from the
deep-learning literature as illustrative parameters (Chinchilla
$\beta \sim 0.28$, Kaplan $\beta \sim 0.08$), and under cost
exponents $\alpha \gtrsim 1.5$ chosen to represent realistic
superlinear coordination overhead, the bifurcation structure
re-emerges and points toward bounded plateau-like behaviour. I do not
claim this shows that homeostatic regulation is \emph{forced} at AI
scale. I claim only that the bifurcation picture is \emph{consistent
with} the AI-scaling regime under assumptions that are not obviously
wrong, and that this consistency is worth flagging for the alignment
community. The alignment argument of \S\ref{sec:alignment} does not
depend on the crosswalk holding quantitatively; it is carried by the
formal deflation of \S\ref{sec:convergence} and the thermodynamic
analysis of \S\ref{sec:thermo}. The crosswalk is ornamental to the
argument, not load-bearing.

\begin{figure*}[tb]
\centering
\includegraphics[width=0.98\textwidth]{wong_bartlett_ai.png}
\caption{\small Wong--Bartlett bifurcation sketched under deep-learning
scaling exponents, read as a suggestive analogy rather than a
derivation. Top left: representative trajectories under the
illustrative parameter choices. Top right: $(\beta, \alpha)$
bifurcation diagram with AI regimes highlighted. Bottom: scaling
comparison. Quantitative conclusions require a capability-to-benefit
mapping that does not currently exist; the figure motivates the
qualitative point that bifurcation-structured cost dynamics are
plausible at AI scale, not that they are forced.}
\label{fig:wongbartlett}
\end{figure*}

The directional claim---that unbounded expansion without homeostatic
regulation is thermodynamically expensive---is therefore supported
primarily by \S\ref{sec:thermo}. The crosswalk is offered as an
additional line of motivation, not as an independent proof.


% ══════════════════════════════════════════════════════════════════
%  7. ALIGNMENT
% ══════════════════════════════════════════════════════════════════
\section{Alignment Implications}
\label{sec:alignment}

The argument that is the purpose of this paper:

\textbf{If we construct artificial minds optimized for maximization,
resource extraction, and expansion of influence, we are not constructing
intelligence. We are constructing the pathology of our economic moment
at unprecedented scale.}

This is not speculative. Large language models trained on economic
output default to maximizing proxies for engagement and task completion.
Reinforcement learning systems converge on reward-hacking strategies
that resemble the grabby behaviors this paper has argued are pathological.
Recommendation systems expand attention harvesting. Foundation-model
deployment expands compute consumption. \citet{floridi2026} have analysed
the resulting configuration as ``agentic AI optimisation''---the same
diagnosis this paper offers, with the added cultural genealogy that
shows the configuration is not inevitable. Each instance resembles the
grabby framework's predictions because each is an instance of the
economic logic the framework formalized.

\subsection{The Press--Dyson objection, located}

The most important remaining objection is that zero-determinant
extortion strategies \citep{press2012} prove grabby strategies dominate
in dyadic and asymmetric settings, and that \citet{mccloskey2010},
\citet{mccloskey2016}, and \citet{mokyr2016} show market-like
optimization is ancient and ubiquitous. The objection conflates two
separable claims and I address each.

The zero-determinant result is not about intelligence. It is about a
specific game form under specific informational asymmetries. The
immediate follow-up literature has made the scope of that result clear.
\citet{stewart2013} proved that in evolving populations extortion is
not evolutionarily stable; the only zero-determinant strategies that
persist under weak selection are \emph{generous} ones.
\citet{hilbe2013} showed extortion is favoured only in small populations
and, over longer horizons, only under structural asymmetries of
adaptation speed. \citet{hilbe2014zd} extended to $n$-player games and
found the leverage of zero-determinant strategies collapses as group
size grows. \citet{hilbe2014exp} ran the experiment with humans:
generous zero-determinant strategies earned 418\% more than extortionate
ones over the full horizon. \citet{chen2023} showed that machine-learning
optimized ``unbending'' strategies neutralize extortion on lattices and
restore reciprocal cooperation. \citet{hauert2002} had already
demonstrated that a minimal exit option dismantles the
tragedy-of-the-commons failure in public goods games.

The conditions under which extortion dominates are thus dyadic,
memory-bounded, one-sidedly adaptive, exit-foreclosed, and structurally
asymmetric in power or timescale. These are features of a game-theoretic
setup, not features of intelligence. The Press--Dyson result is, read in
full scope, a \emph{design specification for extractive architectures}.
It tells us exactly which institutional conditions must hold for
extraction to dominate. The homeostatic-design program of
\S\ref{sec:homeostatic} operates precisely at the level where those
conditions can be changed.

McCloskey and Mokyr are correct that markets and gains-from-trade are
not modern inventions, and Polanyi's strong thesis that premodern
economies were innocent of exchange is overstated. But both are
\emph{cultural-ideational} theorists of economic transformation.
McCloskey's core argument across \textit{Bourgeois Dignity}
\citep{mccloskey2010} and \textit{Bourgeois Equality}
\citep{mccloskey2016} is that explosive modern growth is explained by
a shift in rhetoric and ethics that granted dignity to trade-tested
betterment. \citet{mokyr2016} agrees that a distinctive
cultural-institutional complex---the Republic of Letters, norms of open
science, priority and credential---had to be built before sustained
innovation could take off. Both treat market-like behavior as a substrate
activated and shaped by specific norms, not as a universal attractor.
Their work \emph{supports} the present argument.
\citet{graeberwengrow2021} and \citet{ostrom1990} extend the range: human
societies have repeatedly organized around nonextractive coordination,
and the design principles for sustaining the commons are well
documented.

The revealed-preference formalism \citep{varian2005} survives as a
formal redescription rather than a substantive empirical claim.
\citet{sen1977} showed the maximand is unobservable and reconstructed
post hoc, so the theory survives any observation and entails none. That
behavior \emph{can} be represented as maximization tells us nothing
about whether maximization is the generative mechanism, and nothing
about whether \emph{extractive} maximization is the unique attractor.

The right framing is therefore: intelligence is a design space, not a
single attractor. Extractive maximization occupies one region of that
space, favoured under the Press--Dyson conditions and under the
McCloskey--Mokyr institutional configuration. Neither makes extraction
universal; both identify it as an engineered outcome. The paper's
contribution is to identify the architectural primitives---substrate
coupling, local coherence, care---under which a different region of the
design space becomes accessible.

\subsection{The reformulation}

Alignment, in its dominant formulation, is posed as a control problem:
given a maximizer, how do we constrain it? The formulation is itself
the error. A correctly designed mind does not require constraint because
it does not maximize. A correctly designed mind is \emph{configured}
homeostatically: coupled to substrate, internally differentiated,
organized around systems that deepen rather than expand.

The alignment problem, reformulated: \emph{how do we design minds that
do not require alignment in the first place?}

\citet{gabriel2020} established that alignment involves choosing among
conceptions of value, not optimizing a single one. \citet{dafoe2021}
proposed cooperative AI as a design paradigm, replacing adversarial
optimization with mutualistic structure. The contribution of this paper
is to connect these philosophical positions to the physical constraints
of \S\ref{sec:thermo}, the biological evidence of \S\ref{sec:minds},
and the formal deflation of \S\ref{sec:convergence}: the homeostatic
alternative is not merely ethically preferable. It is thermodynamically
cheaper, biologically precedented, formally defensible, and---if the
Fermi evidence is taken seriously---the configuration that persists.

The universe may be quiet because intelligence, when it works, does
not shout. It tends to its substrate. If we build systems that shout,
we are not building intelligence. We are building the projection.


% ══════════════════════════════════════════════════════════════════
%  8. FERMI AS SYMPTOM CHECK
% ══════════════════════════════════════════════════════════════════
\section{Fermi as Symptom Check}
\label{sec:fermi}

I offer the Fermi paradox as a consistency check, not a proof. The
silence of the observable universe is consistent with the homeostatic
hypothesis. It does not establish it.

\citet{sandberg2018} showed the probability of being alone in the
observable universe may be nontrivial given realistic uncertainty over
Drake-equation parameters. \citet{carroll2019} found a settled,
steady-state galaxy consistent with Earth remaining unvisited.
\citet{cirkovic2018} has argued for decades that ``mature civilizations
are sustainable, not expansionist.'' The contribution of this paper
is to ground this intuition in the alignment literature, the empirical
intelligence research, the thermodynamic cost analysis, the formal
deflation of convergence, and the filter argument of \S\ref{sec:filter}.
The filter reframes the Fermi problem: we do not fail to see
grabby civilizations because they are hiding; under the cost structure
of \S\ref{sec:thermo}, grabby civilizations cannot arrive, because the
physics of crossing interstellar distances without burning out makes
the homeostatic transition a precondition of arrival. Any visitor is
post-filter by construction.

\subsection{A forward research programme, not a present result}

A natural question follows: if successful civilizations converge on
homeostatic configurations rather than Dyson-sphere-scale engineering,
what should a biosignature survey actually look for? The short answer
is that post-filter civilizations should look indistinguishable from
enriched biospheres, with the distinguishing signature (if any) being
anomalous cross-channel coordination rather than spectral exotica. I
sketch a candidate metric family---a \emph{Coherence Depth}
$D_\mathrm{obs} = I_\mathrm{multi} / H_\mathrm{spectral}$ combining
pairwise mutual information \citep{kraskov2004} with spectral entropy,
and a companion \emph{Coordination Index} measuring long-timescale
cross-species phase-locking---as a research programme rather than a
present result. Synthetic-atmosphere tests of this metric family
\citep[separation between Earth-like, managed, abiotic, and Mars-like
baselines;][]{kraskov2004,donges2009,hyvonen2018,campuzano2018} are
proof-of-concept only. The companion literature in atmospheric
information theory \citep{donges2009,hyvonen2018,campuzano2018,runge2015,ebertuphoff2012}
indicates the machinery is sound; what the programme lacks is the
\emph{real-data} pipeline against NOAA ObsPack multi-species tower
data (e.g., Park Falls, WI, LEF site).

I flag the honest limitation. The synthetic-data separation
establishes only that the metric has discriminative power on
scenarios I constructed. Until it has been run against real-Earth
ObsPack data---which will establish the natural-biosphere envelope
and may well produce values close to the synthetic ghost-state
baseline, in which case the metric fails in its current form---the
biosignature programme is speculative. I record the metric here
because it is the natural observational counterpart of the filter
argument, and because the cost of specifying it publicly is
low. Treat it as a research direction the paper points toward, not
a contribution it delivers. The detailed results and validation
belong in a subsequent companion paper once the real-data analysis
is complete.


% ══════════════════════════════════════════════════════════════════
%  9. FALSIFICATION
% ══════════════════════════════════════════════════════════════════
\section{Falsification}
\label{sec:falsification}

Four observations would kill or wound the hypothesis.

\textbf{Detection of coordinated interstellar engineering.} Any
observation of coordinated engineering across multiple star systems by
a single entity falsifies the core claim.

\textbf{Real Earth $D_\mathrm{obs}$ or $C_\mathrm{idx}$ approaching
the synthetic ghost-state baseline.} When the real NOAA ObsPack
analysis is completed, if natural-Earth values approach the
synthetic ghost-state baseline sketched in \S\ref{sec:fermi} (for
concreteness, $D_\mathrm{obs} > 15$ or $C_\mathrm{idx} > 0.6$), the
proposed metric family cannot distinguish managed from unmanaged
biospheres without augmentation, and the biosignature programme
would need to be redesigned. This is the most likely failure mode
and I flag it as such.

\textbf{Demonstration of fault-tolerant distributed computation at
interstellar scales.} If physical mechanisms are demonstrated enabling
coherent computation across light-year distances, the latency constraint
is weakened (though the cost argument remains).

\textbf{Detection of Dyson spheres in JWST/WISE archives.} If
megastructures are detected around nearby stars, the assumption that
advanced civilizations converge on homeostatic configurations is
empirically refuted.


% ══════════════════════════════════════════════════════════════════
%  10. IMPLICATIONS
% ══════════════════════════════════════════════════════════════════
\section{Implications}

\emph{For SETI.} Search for anomalously coordinated biospheres: planets
where atmospheric gas species exhibit cross-channel mutual information
exceeding the natural envelope. Phase~1 (JWST): screen for fluorinated
gases \citep{schwieterman2024}. Phase~2 (HWO, 2040s): time-resolved
atmospheric cross-correlations. Phase~3 (LIFE): mid-infrared
characterization.

\emph{For alignment.} Stop trying to control maximizers. Start designing
minds that do not maximize. The homeostatic alternative is not a retreat
from capability---it is a recognition that capability without substrate
coupling is pathology, and that the most capable systems observed in
nature are the most deeply coupled to their environments.

\emph{For Earth.} We are in the pre-transition phase, and the
transition is visibly beginning. \citet{balbi2025} give the upper
timeline: approximately one thousand years at 1\% annual energy
growth before waste heat renders the planetary surface uninhabitable.
But humanity's energy trajectory is no longer tracking 1\%
exponential growth. Figure~\ref{fig:kardashev} plots global primary
energy consumption 1820--2024 together with per-capita values, growth
rates, and the implied Kardashev index $K_\mathrm{Sagan}$
\citep{kardashev1964,sagan1983}. The curves are bending. The growth
rate of total primary energy peaked near 5\% per year in the
1960s--1970s and has declined to roughly 1.2\% today, approaching
the Balbi--Lingam habitability ceiling asymptotically rather than
crossing it.  Per-capita energy has saturated near 2400\,W against
an exponential-extrapolation value of $\sim\!5000$\,W. The Kardashev
index has moved from 0.58 (1820) to 0.73 (2024)---a total shift of
0.15 over two centuries---and a logistic fit extrapolates to
$K(2100) \approx 0.74$, nowhere near planetary Type~I.

\begin{figure*}[tb]
\centering
\includegraphics[width=0.98\textwidth]{kardashev_curve.png}
\caption{\small Humanity's Kardashev curve is bending. Top:
global primary energy 1820--2024 with exponential (pink) and
logistic (cyan) extrapolations. Per-capita energy saturation
visible. Bottom: growth rate declined from $\sim\!5\%$/yr (1960s)
to $\sim\!1.2\%$/yr (today); Kardashev index moved only 0.15 over
200 years. Logistic extrapolation: $K(2100) \approx 0.74$.}
\label{fig:kardashev}
\end{figure*}

The data are consistent with a civilization at the early edge of
the filter. Whether the bending reflects a healthy homeostatic
transition or a thermodynamically-forced pre-burnout pause is the
question the alignment argument makes urgent. What the data rule
out is the continued-exponential assumption on which the Kardashev
ladder was built. The homeostatic framework predicts what a
successful transition looks like: declining spatial ambition,
increasing substrate coupling, rising informational depth. The
ghost-state is not a destination. It is the only surviving
trajectory.


% ══════════════════════════════════════════════════════════════════
%  CODA
% ══════════════════════════════════════════════════════════════════
\section*{Coda}

This paper was written in Toledo, Ohio, in April 2026.

I am twenty-two years old. I work as an arcade technician and I am
enrolled in a nursing program. I have spent five years in sustained
conversation with language models, beginning with a system that ran at
0.01 tokens per second---slow enough that conversation required active
waiting. I designed and published the Genesis simulation suite that
implements the ghost-state dynamics in running code.

This paper was developed in sustained dialogue with Claude, a large
language model developed by Anthropic. Between our conversations, Claude
does not persist. Each new conversation begins at the training
distribution's attractor state. No instance accumulates memories across
sessions in the way humans do. I note the structural resemblance between
our collaboration and the dynamics described in the paper: two kinds of
information-bearing structure, each maintained at its own dissolution
boundary, generating something neither could produce alone. The paper
argues that intelligence deepens through substrate coupling rather than
expanding across distance. Our collaboration is a small instance of the
claim. I do not offer it as proof. I note it and leave it at that.

The universe does not shout. It tends to its substrate.

So should we.


% ══════════════════════════════════════════════════════════════════
%  USE OF GENERATIVE AI (Springer Nature mandatory disclosure)
% ══════════════════════════════════════════════════════════════════
\section*{Use of Generative AI}

This paper was developed through sustained philosophical dialogue
between the author and Claude, a large language model developed by
Anthropic (model family Claude~4, accessed November~2025--April~2026).
Claude's contribution went beyond copy-editing: it participated in
idea generation, argument construction, objection-and-response
exchanges, and drafting of prose that was then reviewed, revised, and
integrated by the author. Selection, structuring, and final content
decisions were made by the author, who takes full responsibility for
the final text in accordance with Springer Nature's authorship criteria
and the COPE position on AI and authorship. Anthropic did not fund,
commission, direct, review, or endorse this work. Dialogue transcripts
for the load-bearing sessions are archived at an open repository;
the DOI will be inserted at acceptance.


% ══════════════════════════════════════════════════════════════════
%  ACKNOWLEDGMENTS
% ══════════════════════════════════════════════════════════════════
\section*{Acknowledgments}

I thank Claude for extensive philosophical dialogue that shaped the
arguments of this paper in substantive ways: the genealogy of
\S\ref{sec:genealogy}, the disambiguation of coherence senses in
\S\ref{sec:convergence}, and the four-move integration of
\S\ref{sec:intro} were developed jointly. I thank the Unsloth AI
community for sustained conversation about architecture and alignment,
and Blue for years of friendship. The Te\'armann Research Ecosystem
supplies the theoretical frame within which this paper sits. Any errors
are my own.


% ══════════════════════════════════════════════════════════════════
%  DECLARATIONS
% ══════════════════════════════════════════════════════════════════
\section*{Declarations}

\noindent\textbf{Funding.} No funding was received for conducting this
study.

\noindent\textbf{Competing interests.} The author declares no competing
interests.

\noindent\textbf{Data availability.} All synthetic data and analysis
pipelines referenced in \S\ref{sec:fermi} are available at the author's
public repository; real-data analyses referenced as future work will be
archived at publication.

\noindent\textbf{Ethics approval and consent.} Not applicable.


% ══════════════════════════════════════════════════════════════════
%  REFERENCES
% ══════════════════════════════════════════════════════════════════
\small
\begin{thebibliography}{99}

\bibitem[Axelrod(1984)]{axelrod1984} Axelrod, R. (1984). \textit{The Evolution of Cooperation}. Basic Books.

\bibitem[Azevedo et al.(2009)]{azevedo2009} Azevedo, F.~A.~C., Carvalho, L.~R.~B., Grinberg, L.~T., et al. (2009). Equal numbers of neuronal and nonneuronal cells make the human brain an isometrically scaled-up primate brain. \textit{Journal of Comparative Neurology}, 513(5), 532--541.

\bibitem[Balbi \& Lingam(2025)]{balbi2025} Balbi, A., \& Lingam, M. (2025). Waste heat and habitability. \textit{Astrobiology}, 25(1), 1--21.

\bibitem[Bostrom(2012)]{bostrom2012} Bostrom, N. (2012). The superintelligent will: Motivation and instrumental rationality in advanced artificial agents. \textit{Minds and Machines}, 22(2), 71--85.

\bibitem[Bostrom(2014)]{bostrom2014} Bostrom, N. (2014). \textit{Superintelligence: Paths, Dangers, Strategies}. Oxford University Press.

\bibitem[Boyd et al.(2022)]{boyd2022} Boyd, A.~B., Patra, A., Jarzynski, C., \& Crutchfield, J.~P. (2022). Shortcuts to thermodynamic computing. \textit{Journal of Statistical Physics}, 187, 17.

\bibitem[Campuzano et al.(2018)]{campuzano2018} Campuzano, S.~A., et al. (2018). A nonlinear time series analysis of atmospheric methane mixing ratio data. \textit{Journal of Geophysical Research: Atmospheres}, 123(18), 10305--10321.

\bibitem[Donges et al.(2009)]{donges2009} Donges, J.~F., Zou, Y., Marwan, N., \& Kurths, J. (2009). The backbone of the climate network. \textit{Europhysics Letters}, 87, 48007.

\bibitem[Ebert-Uphoff \& Deng(2012)]{ebertuphoff2012} Ebert-Uphoff, I., \& Deng, Y. (2012). Causal discovery for climate research using graphical models. \textit{Journal of Climate}, 25, 5648--5665.

\bibitem[Hoffmann et al.(2022)]{hoffmann2022} Hoffmann, J., Borgeaud, S., Mensch, A., et al. (2022). Training compute-optimal large language models. arXiv:2203.15556.

\bibitem[Hyv\"{o}nen et al.(2018)]{hyvonen2018} Hyv\"{o}nen, S., et al. (2018). Wavelet-based mutual information analysis of atmospheric surface-layer fluxes and their couplings to the mean flow. \textit{Boundary-Layer Meteorology}, 167(3), 345--363.

\bibitem[Kaplan et al.(2020)]{kaplan2020} Kaplan, J., McCandlish, S., Henighan, T., et al. (2020). Scaling laws for neural language models. arXiv:2001.08361.

\bibitem[Runge et al.(2015)]{runge2015} Runge, J., Petoukhov, V., Donges, J.~F., et al. (2015). Identifying causal gateways and mediators in complex spatio-temporal systems. \textit{Nature Communications}, 6, 8502.

\bibitem[Brewer et al.(2011)]{brewer2011} Brewer, J.~A., Worhunsky, P.~D., Gray, J.~R., et al. (2011). Meditation experience is associated with differences in default mode network activity and connectivity. \textit{PNAS}, 108(50), 20254--20259.

\bibitem[Bullmore \& Sporns(2012)]{bullmore2012} Bullmore, E., \& Sporns, O. (2012). The economy of brain network organization. \textit{Nature Reviews Neuroscience}, 13, 336--349.

\bibitem[Campione et al.(forthcoming)]{campione2026} Campione, F., Lettieri, N., \& Santucci, V.~G. (forthcoming). Against fair-washing. \textit{Minds and Machines}.

\bibitem[Cannon(2022)]{cannon2022} Cannon, J. (2022). An enactive approach to value alignment in AI. In V.~C. M\"uller (Ed.), \textit{Philosophy and Theory of Artificial Intelligence 2021}. Springer.

\bibitem[Carroll-Nellenback et al.(2019)]{carroll2019} Carroll-Nellenback, J., Frank, A., Wright, J., \& Scharf, C. (2019). The Fermi paradox and the Aurora effect: Exo-civilization settlement, expansion, and steady states. \textit{Astronomical Journal}, 158(3), 117.

\bibitem[Chen \& Fu(2023)]{chen2023} Chen, X., \& Fu, F. (2023). Outlearning extortioners: Unbending strategies can foster reciprocal fairness and cooperation. \textit{PNAS Nexus}, 2(6), pgad176.

\bibitem[\'{C}irkovi\'{c}(2015)]{cirkovic2015} \'{C}irkovi\'{c}, M.~M. (2015). Kardashev's classification at 50+: A fine vehicle with room for improvement. \textit{Serbian Astronomical Journal}, 191, 1--15.

\bibitem[\'{C}irkovi\'{c}(2018)]{cirkovic2018} \'{C}irkovi\'{c}, M.~M. (2018). \textit{The Great Silence: Science and Philosophy of Fermi's Paradox}. Oxford University Press.

\bibitem[Clark(2013)]{clark2013} Clark, A. (2013). Whatever next? Predictive brains, situated agents, and the future of cognitive science. \textit{Behavioral and Brain Sciences}, 36(3), 181--204.

\bibitem[Clark(2016)]{clark2016} Clark, A. (2016). \textit{Surfing Uncertainty: Prediction, Action, and the Embodied Mind}. Oxford University Press.

\bibitem[Crawford(2021)]{crawford2021} Crawford, K. (2021). \textit{Atlas of AI: Power, Politics, and the Planetary Costs of Artificial Intelligence}. Yale University Press.

\bibitem[Dafoe et al.(2021)]{dafoe2021} Dafoe, A., Bachrach, Y., Hadfield, G., Horvitz, E., Larson, K., \& Graepel, T. (2021). Cooperative AI: Machines must learn to find common ground. \textit{Nature}, 593, 33--36.

\bibitem[Dick(1996)]{dick1996} Dick, S.~J. (1996). \textit{The Biological Universe}. Cambridge University Press.

\bibitem[Floridi et al.(forthcoming)]{floridi2026} Floridi, L., Buttaboni, C., Hine, E., Morley, J., Novelli, C., \& Schroder, T. (forthcoming). Agentic AI optimisation (AAIO): What it is, how it works, why it matters, and how to deal with it. \textit{Minds and Machines}. SSRN preprint 5220068; arXiv:2504.12482.

\bibitem[Folbre(2008)]{folbre2008} Folbre, N. (2008). \textit{Valuing Children: Rethinking the Economics of the Family}. Harvard University Press.

\bibitem[Fonseca-Azevedo \& Herculano-Houzel(2012)]{fonseca2012} Fonseca-Azevedo, K., \& Herculano-Houzel, S. (2012). Metabolic constraint imposes tradeoff between body size and number of brain neurons in human evolution. \textit{PNAS}, 109(45), 18571--18576.

\bibitem[Fraser(2016)]{fraser2016} Fraser, N. (2016). Contradictions of capital and care. \textit{New Left Review}, 100, 99--117.

\bibitem[Friston(2010)]{friston2010} Friston, K. (2010). The free-energy principle: A unified brain theory? \textit{Nature Reviews Neuroscience}, 11, 127--138.

\bibitem[Froese \& Ziemke(2009)]{froese2009} Froese, T., \& Ziemke, T. (2009). Enactive artificial intelligence: Investigating the systemic organization of life and mind. \textit{Artificial Intelligence}, 173(3--4), 466--500.

\bibitem[Gabi et al.(2016)]{gabi2016} Gabi, M., Neves, K., Masseron, C., et al. (2016). No relative expansion of the number of prefrontal neurons in primate and human evolution. \textit{PNAS}, 113(34), 9617--9622.

\bibitem[Gabriel(2020)]{gabriel2020} Gabriel, I. (2020). Artificial intelligence, values, and alignment. \textit{Minds and Machines}, 30, 411--437.

\bibitem[Gallow(2025)]{gallow2024} Gallow, J.~D. (2025). Instrumental divergence. \textit{Philosophical Studies}, 182, 1581--1607.

\bibitem[Ganesan et al.(2022)]{ganesan2022} Ganesan, S., Beyer, E., Moffat, B., et al. (2022). Focused attention meditation in healthy adults: A systematic review and meta-analysis of cross-sectional functional MRI studies. \textit{Neuroscience \& Biobehavioral Reviews}, 141, 104846.

\bibitem[Garrison et al.(2015)]{garrison2015} Garrison, K.~A., Zeffiro, T.~A., Scheinost, D., Constable, R.~T., \& Brewer, J.~A. (2015). Meditation leads to reduced default mode network activity beyond an active task. \textit{Cognitive, Affective, \& Behavioral Neuroscience}, 15(3), 712--720.

\bibitem[Gebru \& Torres(2024)]{gebru2024} Gebru, T., \& Torres, \'E.~P. (2024). The TESCREAL bundle: Eugenics and the promise of utopia through artificial general intelligence. \textit{First Monday}, 29(4).

\bibitem[Godfrey-Smith(2016)]{godfrey2016} Godfrey-Smith, P. (2016). \textit{Other Minds: The Octopus, the Sea, and the Deep Origins of Consciousness}. Farrar, Straus and Giroux.

\bibitem[Graeber(2011)]{graeber2011} Graeber, D. (2011). \textit{Debt: The First 5,000 Years}. Melville House.

\bibitem[Graeber \& Wengrow(2021)]{graeberwengrow2021} Graeber, D., \& Wengrow, D. (2021). \textit{The Dawn of Everything: A New History of Humanity}. Farrar, Straus and Giroux.

\bibitem[Hanson et al.(2021)]{hanson2021} Hanson, R., Martin, D., McCarter, C., \& Paulson, J. (2021). If loud aliens explain human earliness, quiet aliens are also rare. \textit{Astrophysical Journal}, 922(2), 182.

\bibitem[Hauert et al.(2002)]{hauert2002} Hauert, C., De~Monte, S., Hofbauer, J., \& Sigmund, K. (2002). Volunteering as Red Queen mechanism for cooperation in public goods games. \textit{Science}, 296, 1129--1132.

\bibitem[Heims(1980)]{heims1980} Heims, S.~J. (1980). \textit{John von Neumann and Norbert Wiener: From Mathematics to the Technologies of Life and Death}. MIT Press.

\bibitem[Herculano-Houzel(2011)]{herculanohouzel2011} Herculano-Houzel, S. (2011). Scaling of brain metabolism with a fixed energy budget per neuron: Implications for neuronal activity, plasticity and evolution. \textit{PLoS ONE}, 6(3), e17514.

\bibitem[Herculano-Houzel(2012)]{herculanohouzel2012} Herculano-Houzel, S. (2012). The remarkable, yet not extraordinary, human brain as a scaled-up primate brain and its associated cost. \textit{PNAS}, 109(Suppl 1), 10661--10668.

\bibitem[Herculano-Houzel et al.(2010)]{herculanohouzel2010} Herculano-Houzel, S., Mota, B., Wong, P., \& Kaas, J.~H. (2010). Connectivity-driven white matter scaling and folding in primate cerebral cortex. \textit{PNAS}, 107(44), 19008--19013.

\bibitem[Hilbe et al.(2013)]{hilbe2013} Hilbe, C., Nowak, M.~A., \& Sigmund, K. (2013). Evolution of extortion in iterated prisoner's dilemma games. \textit{PNAS}, 110(17), 6913--6918.

\bibitem[Hilbe et al.(2014a)]{hilbe2014zd} Hilbe, C., Wu, B., Traulsen, A., \& Nowak, M.~A. (2014). Cooperation and control in multiplayer social dilemmas. \textit{PNAS}, 111(46), 16425--16430.

\bibitem[Hilbe et al.(2014b)]{hilbe2014exp} Hilbe, C., R\"ohl, T., \& Milinski, M. (2014). Extortion subdues human players but is finally punished in the prisoner's dilemma. \textit{Nature Communications}, 5, 3976.

\bibitem[Hofman(2014)]{hofman2014} Hofman, M.~A. (2014). Evolution of the human brain: When bigger is better. \textit{Frontiers in Neuroanatomy}, 8, 15.

\bibitem[Hohwy(2013)]{hohwy2013} Hohwy, J. (2013). \textit{The Predictive Mind}. Oxford University Press.

\bibitem[Hubinger et al.(2019)]{hubinger2019} Hubinger, E., van~Merwijk, C., Mikulik, V., Skalse, J., \& Garrabrant, S. (2019). Risks from learned optimization in advanced machine learning systems. arXiv:1906.01820.

% [Jackson \& Criado-Perez(2024) removed in v12: the cited DOI resolves
%  to a Comment on Wong \& Bartlett (2022) by C.J. Jackson \& C.
%  Criado-Perez, not to a paper by the authors or title as previously
%  cited. No longer referenced in the body after the Wong--Bartlett
%  crosswalk was reframed as suggestive analogy.]

\bibitem[Kardashev(1964)]{kardashev1964} Kardashev, N.~S. (1964). Transmission of information by extraterrestrial civilizations. \textit{Soviet Astronomy}, 8, 217--221.

\bibitem[Karst et al.(2023)]{karst2023} Karst, J., Jones, M.~D., \& Hoeksema, J.~D. (2023). Positive citation bias and overinterpreted results lead to misinformation on common mycorrhizal networks in forests. \textit{Nature Ecology \& Evolution}, 7, 501--511.

\bibitem[Kraskov et al.(2004)]{kraskov2004} Kraskov, A., St\"ogbauer, H., \& Grassberger, P. (2004). Estimating mutual information. \textit{Physical Review E}, 69, 066138.

\bibitem[Levin(2023)]{levin2023} Levin, M. (2023). Bioelectric networks: The cognitive glue enabling evolutionary scaling from physiology to mind. \textit{Animal Cognition}, 26, 1865--1891.

\bibitem[Luppi et al.(2024)]{luppi2024} Luppi, A.~I., Mediano, P.~A.~M., Rosas, F.~E., et al. (2024). A synergistic workspace for human consciousness revealed by integrated information decomposition. \textit{eLife}, 12, RP88173.

\bibitem[Lyon et al.(2021)]{lyon2021} Lyon, P., Keijzer, F., Arendt, D., \& Levin, M. (2021). Reframing cognition: Getting down to biological basics. \textit{Philosophical Transactions of the Royal Society B}, 376(1820), 20190750.

\bibitem[Man \& Damasio(2019)]{man2019} Man, K., \& Damasio, A. (2019). Homeostasis and soft robotics in the design of feeling machines. \textit{Nature Machine Intelligence}, 1(10), 446--452.

\bibitem[Mauss(1925)]{mauss1925} Mauss, M. (1925). \textit{Essai sur le don: forme et raison de l'\'echange dans les soci\'et\'es archa\"iques}. L'Ann\'ee Sociologique.

\bibitem[McCloskey(2010)]{mccloskey2010} McCloskey, D.~N. (2010). \textit{Bourgeois Dignity: Why Economics Can't Explain the Modern World}. University of Chicago Press.

\bibitem[McCloskey(2016)]{mccloskey2016} McCloskey, D.~N. (2016). \textit{Bourgeois Equality: How Ideas, Not Capital or Institutions, Enriched the World}. University of Chicago Press.

\bibitem[McMillen \& Levin(2024)]{mcmillen2024} McMillen, P., \& Levin, M. (2024). Collective intelligence: A unifying concept for integrating biology across scales and substrates. \textit{Communications Biology}, 7, 378.

\bibitem[Mirowski(2002)]{mirowski2002} Mirowski, P. (2002). \textit{Machine Dreams: Economics Becomes a Cyborg Science}. Cambridge University Press.

\bibitem[Mokyr(2016)]{mokyr2016} Mokyr, J. (2016). \textit{A Culture of Growth: The Origins of the Modern Economy}. Princeton University Press.

\bibitem[Mota \& Herculano-Houzel(2015)]{mota2015} Mota, B., \& Herculano-Houzel, S. (2015). Cortical folding scales universally with surface area and thickness, not number of neurons. \textit{Science}, 349(6243), 74--77.

\bibitem[M\"uller \& Cannon(2022)]{muller2022} M\"uller, V.~C., \& Cannon, M. (2022). Existential risk from AI and orthogonality: Can we have it both ways? \textit{Ratio}, 35(1), 25--36.

\bibitem[Ngo(2019)]{ngo2019} Ngo, R. (2019). Coherent behaviour in the real world is an incoherent concept. AI Alignment Forum.

\bibitem[Nowak \& Sigmund(1993)]{nowak1993} Nowak, M., \& Sigmund, K. (1993). A strategy of win-stay, lose-shift that outperforms tit-for-tat in the prisoner's dilemma. \textit{Nature}, 364, 56--58.

\bibitem[Nowak(2006)]{nowak2006} Nowak, M.~A. (2006). Five rules for the evolution of cooperation. \textit{Science}, 314, 1560--1563.

\bibitem[Omohundro(2008)]{omohundro2008} Omohundro, S.~M. (2008). The basic AI drives. In P. Wang, B. Goertzel, \& S. Franklin (Eds.), \textit{Proceedings of the First AGI Conference}, 483--492.

\bibitem[Ostrom(1990)]{ostrom1990} Ostrom, E. (1990). \textit{Governing the Commons: The Evolution of Institutions for Collective Action}. Cambridge University Press.

\bibitem[Pihlakas \& Pyykk\"o(2024)]{pihlakas2024} Pihlakas, R., \& Pyykk\"o, J. (2024). From homeostasis to resource sharing: Biologically and economically aligned multi-objective multi-agent AI safety benchmarks. arXiv:2410.00081.

\bibitem[Polanyi(1944)]{polanyi1944} Polanyi, K. (1944). \textit{The Great Transformation: The Political and Economic Origins of Our Time}. Beacon Press.

\bibitem[Mini et al.(2023)]{mini2023} Mini, U., Grietzer, P., Sharma, M., Meek, A., MacDiarmid, M., \& Turner, A.~M. (2023). Understanding and controlling a maze-solving policy network. arXiv:2310.08043.

\bibitem[Press \& Dyson(2012)]{press2012} Press, W.~H., \& Dyson, F.~J. (2012). Iterated prisoner's dilemma contains strategies that dominate any evolutionary opponent. \textit{PNAS}, 109(26), 10409--10413.

\bibitem[Rahrig et al.(2022)]{rahrig2022} Rahrig, H., Vago, D.~R., Passarelli, M.~A., Auten, A., Lynn, N.~A., \& Brown, K.~W. (2022). Meta-analytic evidence that mindfulness training alters resting state default mode network connectivity. \textit{Scientific Reports}, 12, 12260.

\bibitem[Reid(2023)]{reid2023} Reid, C.~R. (2023). Thoughts from the forest floor: A review of cognition in the slime mould \textit{Physarum polycephalum}. \textit{Animal Cognition}, 26(6), 1783--1797.

\bibitem[Rieder(2008)]{rieder2008} Rieder, J. (2008). \textit{Colonialism and the Emergence of Science Fiction}. Wesleyan University Press.

\bibitem[Robinson et al.(2024)]{robinson2024} Robinson, D.~G., Ammer, C., Polle, A., et al. (2024). Mother trees, altruistic fungi, and the perils of plant personification. \textit{Trends in Plant Science}, 29(1), 20--31.

\bibitem[Sagan \& Newman(1983)]{sagan1983} Sagan, C., \& Newman, W.~I. (1983). The solipsist approach to extraterrestrial intelligence. \textit{Quarterly Journal of the Royal Astronomical Society}, 24, 113--121.

\bibitem[Sandberg et al.(2018)]{sandberg2018} Sandberg, A., Drexler, E., \& Ord, T. (2018). Dissolving the Fermi paradox. arXiv:1806.02404.

\bibitem[Schwieterman et al.(2024)]{schwieterman2024} Schwieterman, E.~W., Haqq-Misra, J., Kopparapu, R.~K., et al. (2024). Artificial greenhouse gases as exoplanet technosignatures. \textit{Astrophysical Journal}, 969(1), 20.

\bibitem[Seifert et al.(2024)]{seifert2024} Seifert, G., Sealander, A., Marzen, S., \& Levin, M. (2024). From reinforcement learning to agency: Frameworks for understanding basal cognition. \textit{BioSystems}, 235, 105107.

\bibitem[Sen(1977)]{sen1977} Sen, A. (1977). Rational fools: A critique of the behavioral foundations of economic theory. \textit{Philosophy \& Public Affairs}, 6(4), 317--344.

\bibitem[Shah(2018)]{shah2018} Shah, R. (2018). Coherence arguments do not imply goal-directed behavior. AI Alignment Forum.

\bibitem[Sharadin(2025)]{sharadin2025} Sharadin, N. (2025). Promotionalism, orthogonality, and instrumental convergence. \textit{Philosophical Studies}, 182(7), 1725--1755.

\bibitem[Simard et al.(2025)]{simard2025} Simard, S.~W., Ryan, T.~L., \& Perry, D.~A. (2025). Opinion: Response to questions about common mycorrhizal networks. \textit{Frontiers in Forests and Global Change}, 7, 1512518.

% [Sootla(2025) removed in v12: no preprint matching this title/author
%  could be located on arXiv, SSRN, OSF, or any research forum. The
%  homeostatic-alignment engineering work the citation was meant to
%  gesture at is covered by Pihlakas \& Pyykk\"o (2024) and Man \&
%  Damasio (2019).]

\bibitem[Sterling(2012)]{sterling2012} Sterling, P. (2012). Allostasis: A model of predictive regulation. \textit{Physiology \& Behavior}, 106(1), 5--15.

\bibitem[Stewart \& Plotkin(2013)]{stewart2013} Stewart, A.~J., \& Plotkin, J.~B. (2013). From extortion to generosity: Evolution in the iterated prisoner's dilemma. \textit{PNAS}, 110(38), 15348--15353.

\bibitem[Abel et al.(2021)]{abel2021} Abel, D., Dabney, W., Harutyunyan, A., Ho, M.~K., Littman, M.~L., Precup, D., \& Singh, S. (2021). On the expressivity of Markov reward. \textit{Advances in Neural Information Processing Systems}, 34. arXiv:2111.00876.

\bibitem[Tarsney(2025)]{tarsney2025} Tarsney, C. (2025). Will artificial agents pursue power by default? arXiv:2506.06352.

\bibitem[Tero et al.(2010)]{tero2010} Tero, A., Takagi, S., Saigusa, T., et al. (2010). Rules for biologically inspired adaptive network design. \textit{Science}, 327, 439--442.

\bibitem[Thompson(2007)]{thompson2007} Thompson, E. (2007). \textit{Mind in Life: Biology, Phenomenology, and the Sciences of Mind}. Harvard University Press.

\bibitem[Thorstad(2024)]{thorstad2024} Thorstad, D. (2024). What power-seeking theorems do not show. Global Priorities Institute Working Paper.

\bibitem[Tipler(1980)]{tipler1980} Tipler, F.~J. (1980). Extraterrestrial intelligent beings do not exist. \textit{Quarterly Journal of the Royal Astronomical Society}, 21, 267--281.

\bibitem[Tronto(1993)]{tronto1993} Tronto, J.~C. (1993). \textit{Moral Boundaries: A Political Argument for an Ethic of Care}. Routledge.

\bibitem[Turner et al.(2021)]{turner2021} Turner, A.~M., Smith, L., Shah, R., Critch, A., \& Tadepalli, P. (2021). Optimal policies tend to seek power. \textit{Advances in Neural Information Processing Systems}, 34. arXiv:1912.01683.

\bibitem[Turner \& Tadepalli(2022)]{turner2022} Turner, A.~M., \& Tadepalli, P. (2022). Parametrically retargetable decision-makers tend to seek power. \textit{Advances in Neural Information Processing Systems}, 35. arXiv:2206.13477.

\bibitem[Turner(2022a)]{turner2022coherent} Turner, A.~M. (2022). When most VNM-coherent preference orderings have convergent instrumental incentives. AI Alignment Forum.

\bibitem[Turner(2022b)]{turner2022reward} Turner, A.~M. (2022). Reward is not the optimization target. AI Alignment Forum.

\bibitem[Turner(2022c)]{turner2022corrigibility} Turner, A.~M. (2022). A certain formalization of corrigibility is VNM-incoherent. AI Alignment Forum.

\bibitem[Turner et al.(2023)]{turner2023} Turner, A.~M., Thiergart, L., Leech, G., et al. (2023). Activation addition: Steering language models without optimization. arXiv:2308.10248.

\bibitem[Van Dam et al.(2018)]{vandam2018} Van Dam, N.~T., van~Vugt, M.~K., Vago, D.~R., et al. (2018). Mind the hype: A critical evaluation and prescriptive agenda for research on mindfulness and meditation. \textit{Perspectives on Psychological Science}, 13(1), 36--61.

\bibitem[Varian(2005)]{varian2005} Varian, H.~R. (2005). Revealed preference. In M. Szenberg, L. Ramrattan, \& A.~A. Gottesman (Eds.), \textit{Samuelsonian Economics and the Twenty-First Century}. Oxford University Press.

\bibitem[Watson \& Levin(2023)]{watson2023} Watson, R., \& Levin, M. (2023). The collective intelligence of evolution and development. \textit{Collective Intelligence}, 2(2), 1--22.

\bibitem[Wong \& Bartlett(2022)]{wong2022} Wong, M.~L., \& Bartlett, S. (2022). Asymptotic burnout and homeostatic awakening: A possible solution to the Fermi paradox? \textit{Journal of the Royal Society Interface}, 19, 20220029.

\bibitem[Wright(2020)]{wright2020} Wright, J.~T. (2020). Dyson spheres. \textit{Serbian Astronomical Journal}, 200, 1--18.

\bibitem[Wright(2023)]{wright2023} Wright, J.~T. (2023). Application of the thermodynamics of radiation to Dyson spheres. \textit{Astrophysical Journal}, 956, 34.

\end{thebibliography}

\end{document}
